\chapter{FX Derivatives}\label{ch:dv:fx-derivatives}

On 11 August 2015 the People's Bank of China changed how it set the renminbi's daily central parity against the dollar.
The parity would now follow the previous day's closing market rate rather than a preset target. In the two days that
followed, the renminbi fell 2.8\% against the dollar, and currencies across Asia fell with it. For a client holding a
target-redemption forward, a structure that pays a little while the rate stays on one side of a strike and loses twice as
fast on the other, a move of that size is the difference between a product that redeems early with a small gain and one
that runs to maturity at a large loss. Currency options are quoted in a highly standardised way, in volatility by delta, and their exotics are
priced from those quotes. This chapter revisits the pricing conventions, builds the
market's own smile-adjustment rule, vanna--volga, prices barriers and target-redemption forwards under models that fit the
same smile, and introduces the model that combines local and stochastic volatility.

\section{Pricing conventions revisited}

One Quant Book 2, chapter 19, set out how currency options are quoted: in volatility, by delta, with the at-the-money
straddle, the risk reversal and the butterfly as the three quoted instruments per expiry. The model under the quotes is the
following.

\begin{definition}[Garman--Kohlhagen model]\label{def:dv:fx-derivatives:gk}
The \emph{Garman--Kohlhagen model}\index{Garman--Kohlhagen model} is Black--Scholes for an exchange rate $S$ (units of the quote
currency per unit of the base currency): the base currency is an asset paying a continuous yield equal to its interest rate $r_f$,
the quote currency's rate $r_d$ discounts, and a call on the base currency is worth
$e^{-r_fT}S\,\Phi(d_1)-e^{-r_dT}K\,\Phi(d_2)$ with $d_{1,2}=\bigl(\ln(F/K)\pm\frac12\sigma^2T\bigr)/(\sigma\sqrt T)$ and
$F=Se^{(r_d-r_f)T}$.
\end{definition}

Three conventions that the model leaves open decide every quoted number: which delta (spot or forward, premium-adjusted or not),
which at-the-money (the delta-neutral straddle for most pairs), and how the butterfly is defined. Book 2's component
\texttt{firm.fxsmile} implements them. This chapter takes a pair whose market is known exactly, so that every method can be checked
against the truth. The pair is a dollar against an emerging currency, with spot 7.00, a quote-currency rate of 3\% and a dollar rate
of 5\%, and its market is a \omterm{def:dv:stochastic-volatility:heston}{Heston model} (chapter 10) with $v_0=0.0036$, $\kappa=1.5$, $\bar v=0.0049$, $\eta=0.30$ and $\rho=+0.5$. The
correlation is positive because in such pairs the dollar's calls, the protection against the emerging currency's fall, are the
expensive side.

\begin{example}[The broker's screen]\label{ex:dv:fx-derivatives:quotes}
At one year the market implies an at-the-money volatility of 5.32\% (delta-neutral straddle strike 6.871, against a forward of 6.861). The
25-delta risk reversal is $+1.89$ volatility points and the butterfly $+0.46$, with the 25-delta strikes at 6.662 and 7.177. At 10 delta
the risk reversal is $+3.77$ and the butterfly $+1.71$ (strikes 6.442 and 7.702).
\end{example}

\section{The vanna--volga method}

Dealers in barrier and touch options needed a rule that prices them from the three quoted instruments without a full smile
model. The rule they adopted hedges the exotic's volatility risk with the three
instruments and charges their cost.

\begin{definition}[Vanna--volga method]\label{def:dv:fx-derivatives:vv}
The \emph{vanna--volga method}\index{vanna--volga method} prices an option as its Black price at the at-the-money volatility plus
the smile cost of a portfolio of the three quoted instruments (25-delta put, at-the-money, 25-delta call) that matches the option's
\omterm{def:dv:greeks-and-the-hedging-pnl:greeks}{vega}, \omterm{def:dv:greeks-and-the-hedging-pnl:greeks}{vanna} and \omterm{def:dv:greeks-and-the-hedging-pnl:greeks}{volga} (chapter 4), all computed at the at-the-money volatility.
\end{definition}

For a vanilla at strike $K$ the weights $w_i$ solve a three-by-three linear system in the pillars' \omterm{def:dv:greeks-and-the-hedging-pnl:greeks}{vega}, \omterm{def:dv:greeks-and-the-hedging-pnl:greeks}{vanna} and \omterm{def:dv:greeks-and-the-hedging-pnl:greeks}{volga}, and the price
is $C_{\mathrm{BS}}(K,\sigma_{\mathrm{ATM}})+\sum_iw_i\bigl(C^{\mathrm{mkt}}_i-C^{\mathrm{BS}}_i\bigr)$. It reproduces the three quotes
exactly, and between them it gives a smooth smile, which is why the method is widely used in currency markets to build a whole smile
from three quotes. Its underlying assumption has been described as a ``flat but stochastic'' implied volatility. Outside the pillars
it extrapolates, and there it can be wrong.

\begin{example}[Vanna--volga against the true smile]\label{ex:dv:fx-derivatives:vvsmile}
Built from the three 25-delta pillars of \cref{ex:dv:fx-derivatives:quotes}, the vanna--volga smile matches the market to a few hundredths
of a point between the pillars. At the 10-delta call strike it gives 7.94\% against the market's 8.91\%, a point low
(\cref{fig:dv:fx-derivatives:smile}). The 10-delta quotes carry information that the three pillars do not.
\end{example}

\begin{omfigure}
\begin{tikzpicture}
\begin{axis}[width=10.5cm,height=5cm,
  xlabel={strike},ylabel={1-year implied volatility (\%)},
  tick label style={font=\small},label style={font=\small},
  xmin=6.35,xmax=7.8,ymin=4,ymax=9.8,grid=major,grid style={gray!25},
  legend columns=3,legend cell align=left,legend style={font=\footnotesize,at={(0.5,-0.3)},anchor=north,draw=none,fill=none,/tikz/every even column/.append style={column sep=8pt}},
  table/col sep=comma]
\addplot[gray,very thick] table[x=k,y=market]{figdata/derivatives/20-fx-derivatives/smile.csv};
\addplot[omThm,very thick,dashed] table[x=k,y=vv]{figdata/derivatives/20-fx-derivatives/smile.csv};
\addplot[black,only marks,mark=*,mark size=2pt] table[x=k,y=vol]{figdata/derivatives/20-fx-derivatives/pillars.csv};
\legend{market (Heston),vanna--volga,the three pillars}
\end{axis}
\end{tikzpicture}

\omcaption{One-year smile of the chapter's pair from 10-delta put to 10-delta call: the market, and the vanna--volga smile built from the
25-delta put, the at-the-money straddle and the 25-delta call. Exact at the pillars, close between them, low in the call wing. Data: the
tutorial.}\label{fig:dv:fx-derivatives:smile}
\end{omfigure}

For a touch or a barrier the same idea is applied to the exotic's own \omterm{def:dv:greeks-and-the-hedging-pnl:greeks}{vanna} and \omterm{def:dv:greeks-and-the-hedging-pnl:greeks}{volga}. The costs are those of the risk reversal (the
market price of \omterm{def:dv:greeks-and-the-hedging-pnl:greeks}{vanna}) and of the butterfly (the market price of \omterm{def:dv:greeks-and-the-hedging-pnl:greeks}{volga}). Market practice then scales the correction by a probability,
commonly the probability of not touching the barrier, because a barrier that is hit early stops needing a smile hedge. The scaling is
a convention, not a derivation, and different desks use different ones. The next section measures what it does.

\section{Barriers and target-redemption forwards}

\begin{example}[One-touches under five prices]\label{ex:dv:fx-derivatives:onetouch}
A one-year one-touch paying 1 at expiry if the rate touches an upper barrier, continuous monitoring (\cref{fig:dv:fx-derivatives:onetouch}).
For a barrier at 7.20, Black at the at-the-money volatility gives 45.5\% of the payout, vanna--volga 34.5\%, \omterm{def:dv:local-volatility:model}{local volatility} 43.3\%,
stochastic-local volatility 42.2\% and Heston 38.9\%. At 7.60 they give 6.0\%, 21.8\%, 14.3\%, 14.0\% and 13.3\%. At 8.00 they give 0.4\%,
6.7\%, 6.5\%, 6.4\% and 6.2\%. The three models fit the same smile and disagree by up to 4.4 points. The at-the-money price is far off,
and vanna--volga's convention overshoots the models at the middle barriers by more than half.
\end{example}

\begin{omfigure}
\begin{tikzpicture}
\begin{axis}[width=10.5cm,height=5.4cm,
  xlabel={upper barrier},ylabel={one-touch price (\% of payout)},
  tick label style={font=\small},label style={font=\small},
  xmin=7.15,xmax=8.05,ymin=0,ymax=48,grid=major,grid style={gray!25},
  legend columns=3,legend cell align=left,legend style={font=\footnotesize,at={(0.5,-0.28)},anchor=north,draw=none,fill=none,/tikz/every even column/.append style={column sep=8pt}},
  table/col sep=comma]
\addplot[gray,thick,dashed,mark=o,mark size=1.6pt] table[x=h,y=bs]{figdata/derivatives/20-fx-derivatives/one_touch.csv};
\addplot[omProp,thick,mark=triangle*,mark size=1.8pt] table[x=h,y=vv]{figdata/derivatives/20-fx-derivatives/one_touch.csv};
\addplot[omDef,very thick,mark=square*,mark size=1.6pt] table[x=h,y=lv]{figdata/derivatives/20-fx-derivatives/one_touch.csv};
\addplot[black,thick,mark=diamond*,mark size=1.8pt] table[x=h,y=slv]{figdata/derivatives/20-fx-derivatives/one_touch.csv};
\addplot[omThm,very thick,mark=*,mark size=1.6pt] table[x=h,y=sv]{figdata/derivatives/20-fx-derivatives/one_touch.csv};
\legend{Black at ATM,vanna--volga,local volatility,{stochastic-local, mixing 0.5},Heston}
\end{axis}
\end{tikzpicture}

\omcaption{One-year upper one-touches on the chapter's pair (spot 7.00) by barrier. The three models fit the same smile; Black at the
at-the-money volatility ignores it; vanna--volga applies the market's rule of thumb. Data: the tutorial.}\label{fig:dv:fx-derivatives:onetouch}
\end{omfigure}

The models' disagreement is the dynamics of chapters 9 to 11 again. \omterm{def:dv:local-volatility:model}{Local volatility} makes the smile a function of the spot, so as the
rate rises toward the barrier it enters the high-volatility call wing and touches more often. Heston keeps part of the volatility in a
variance process, correlated with the rate but not a function of it. A touch depends on the whole path, and so on those dynamics, not
only on the terminal smile.

The currency market's most notorious \omterm{def:dv:the-structured-products-business:product}{structured products} are built on forwards rather than options.

\begin{definition}[Target redemption forward]\label{def:dv:fx-derivatives:tarf}
A \emph{target redemption forward}\index{target redemption forward} (TARF) is a strip of forward contracts at an enhanced strike, one
per fixing date, that terminates once the client's accumulated gains reach a target. Fixings on the unfavourable side of the strike are
settled on a leveraged notional (typically twice), with no limit on the losses.
\end{definition}

\begin{definition}[Accumulator]\label{def:dv:fx-derivatives:accumulator}
An \emph{accumulator}\index{accumulator} is the equity or currency analogue in which the client buys (or sells) a fixed quantity at a
discounted strike on each fixing while a knock-out barrier has not been touched, and twice the quantity when the price is on the wrong
side of the strike.
\end{definition}

Both give the client a better rate than the forward in exchange for a short option position with leverage. The target or the knock-out
caps the upside quickly, and the leverage leaves the downside open. The client has sold volatility, and usually more than he realises.
The product is priced at zero cost by choosing the strike.

\begin{example}[A TARF on the pair]\label{ex:dv:fx-derivatives:tarf}
The client sells one million dollars at each of twelve monthly fixings at a strike $K$: a gain of $K-S$ per dollar when the fixing is below
$K$, a loss of $2(S-K)$ when above, and redemption once the gains reach 0.30 per dollar. Under stochastic-local volatility (mixing 0.5) the
zero-cost strike is 7.312, against a spot of 7.00 and a one-year forward of 6.861. The client is offered a rate 6.6\% better than the
forward. The TARF redeems in 98.1\% of paths, after 1.73 fixings on average, and 95\% of paths end with a gain of at least 126\,000
(\cref{fig:dv:fx-derivatives:tarf}, left). The zero value comes from the other 5\%. If the rate rises by 10\% over the first quarter, the
client's expected loss is 7.12 million per million of monthly notional: 1.01 million in the first quarter and the rest over nine more
fixings at twice the notional. A rise of 5\% costs 0.99 million, while a rise of up to 2.5\% still lets the TARF redeem with the full
0.30 million (\cref{fig:dv:fx-derivatives:tarf}, right).
\end{example}

\begin{omfigure}
\begin{tikzpicture}
\begin{groupplot}[group style={group size=2 by 1,horizontal sep=1.6cm},
  width=6.6cm,height=5cm,
  tick label style={font=\small},label style={font=\small},grid=major,grid style={gray!25},table/col sep=comma]
\nextgroupplot[ybar,bar width=6pt,xlabel={fixings until redemption},ylabel={share of paths (\%)},xmin=0.3,xmax=12.7,ymin=0,ymax=72,
  xtick={1,2,4,6,8,10,12},title={\small how long it lives}]
\addplot[omThm,fill=omThm!40] table[x=n,y=share]{figdata/derivatives/20-fx-derivatives/tarf_fixings.csv};
\nextgroupplot[xlabel={rate move over the first quarter (\%)},ylabel={client's expected P\&L (million)},xmin=-6,xmax=11,ymin=-8,ymax=1,
  title={\small a move against the client}]
\addplot[omDef,very thick,mark=*,mark size=1.6pt] table[x=move,y=pnl]{figdata/derivatives/20-fx-derivatives/tarf_scenarios.csv};
\draw[gray] (axis cs:-6,0) -- (axis cs:11,0);
\end{groupplot}
\end{tikzpicture}

\omcaption{A zero-cost TARF on the chapter's pair (client sells the dollar at 7.312, twice the notional above it, target 0.30). Left: the
number of fixings before redemption: most end at the first or second. Right: the client's expected P\&L, per million of monthly notional,
if the rate moves by a given amount over the first quarter; the fixings of a TARF that fails to redeem are paid at twice the
notional. Data: the tutorial.}\label{fig:dv:fx-derivatives:tarf}
\end{omfigure}

The August 2015 move shows why the product is dangerous. For a client in a TARF of this kind, a sustained depreciation of the emerging
currency is exactly the scenario the zero-cost strike was paid for. The target never fills, and the leveraged fixings accumulate. The
dealer's side is no simpler. A TARF's delta flips as the accumulated gain approaches the target, and its value depends on the smile's
dynamics through the path, like the touches above.

\section{Stochastic-local volatility}

\omterm{def:dv:local-volatility:model}{Local volatility} fits the smile and gets the dynamics wrong in one direction; stochastic volatility gets them partly right and cannot
fit the smile exactly. The natural compromise combines them.

\begin{definition}[Stochastic-local volatility model]\label{def:dv:fx-derivatives:slv}
A \emph{stochastic-local volatility model}\index{stochastic-local volatility model} (SLV) lets the underlying's volatility be the
product of a function of time and spot and a stochastic factor:
\[
\frac{dS_t}{S_t}=(r_d-r_f)\,dt+L(t,S_t)\sqrt{v_t}\,dW^1_t,\qquad dv_t=\kappa(\bar v-v_t)\,dt+\lambda\eta\sqrt{v_t}\,dW^2_t,
\]
with a mixing weight $\lambda\in[0,1]$ scaling the volatility of variance: $\lambda=0$ reduces the model to \omterm{def:dv:local-volatility:model}{local volatility} and $\lambda=1$
with $L(t,S)=1$ to Heston.
\end{definition}

\begin{definition}[Leverage function]\label{def:dv:fx-derivatives:leverage}
The \emph{leverage function}\index{leverage function} $L(t,S)$ of a \omterm{def:dv:fx-derivatives:slv}{stochastic-local volatility model} is the function that makes the model
reprice every vanilla. By the \omterm{def:dv:local-volatility:projection}{Markovian projection} of chapter 9 it satisfies
\[
L(t,S)^2\,\E\bigl[v_t\mid S_t=S\bigr]=\sigma_{\mathrm{loc}}(t,S)^2,
\]
with $\sigma_{\mathrm{loc}}(t,S)$ the market's \omterm{def:dv:local-volatility:model}{local volatility}.
\end{definition}

The equation is circular: $L(t,S)$ depends on the conditional expectation, which depends on the law of the paths, which depends on $L(t,S)$.
Guyon and Henry-Labord\`ere showed that the particle method of McKean resolves it exactly. Simulate a large set of particles forward in time.
At each step estimate $\E[v_t\mid S_t=S]$ from the particles themselves, for instance by sorting them into bins in $S$, and set
$L(t,S)$ from the equation before taking the next step. The calibration costs one simulation. The build does it with 50\,000 particles
and 40 bins.

\begin{example}[The leverage function of the pair]\label{ex:dv:fx-derivatives:leverage}
With mixing $\lambda=0.5$, the calibrated $L(t,S)$ near the forward is 0.76 at three months, 0.73 at six months and 0.79 at one year. It
rises to about 1.4 or 1.5 in both wings (\cref{fig:dv:fx-derivatives:leverage}). All three models reprice the one-year smile: at the 10-delta
put, the 25-delta put, the at-the-money straddle, the 25-delta call and the 10-delta call, stochastic-local volatility gives 5.22\%, 4.90\%,
5.38\%, 6.77\% and 8.96\%, against the market's 5.14\%, 4.84\%, 5.32\%, 6.72\% and 8.91\%. The differences, below a tenth of a point, come
from the simulation's daily Euler steps.
\end{example}

\begin{omfigure}
\begin{tikzpicture}
\begin{axis}[width=10.5cm,height=5cm,
  xlabel={spot},ylabel={leverage $L(t,S)$},
  tick label style={font=\small},label style={font=\small},
  xmin=6.0,xmax=8.6,ymin=0.6,ymax=1.6,grid=major,grid style={gray!25},
  legend columns=3,legend cell align=left,legend style={font=\footnotesize,at={(0.5,-0.3)},anchor=north,draw=none,fill=none,/tikz/every even column/.append style={column sep=8pt}},
  table/col sep=comma]
\addplot[omProp,thick] table[x=s3m,y=l3m]{figdata/derivatives/20-fx-derivatives/leverage_axes.csv};
\addplot[omDef,very thick,dashed] table[x=s6m,y=l6m]{figdata/derivatives/20-fx-derivatives/leverage_axes.csv};
\addplot[omThm,very thick] table[x=s12m,y=l12m]{figdata/derivatives/20-fx-derivatives/leverage_axes.csv};
\legend{3 months,6 months,1 year}
\end{axis}
\end{tikzpicture}

\omcaption{The \omterm{def:dv:fx-derivatives:leverage}{leverage function} of the \omterm{def:dv:fx-derivatives:slv}{stochastic-local volatility model} (mixing 0.5) calibrated by the particle method to the chapter's pair,
at three dates, over the range the particles cover. Below one near the forward, above one in the wings: the local part supplies the smile that
half the volatility of variance leaves out. Data: the tutorial.}\label{fig:dv:fx-derivatives:leverage}
\end{omfigure}

The mixing weight is the model's free parameter. It is not fitted to vanillas, which every $\lambda$ reprices. It is chosen from the
products that depend on dynamics, touches and barriers above all, by calibrating it to their prices when the market quotes them, or to
the \omterm{def:dv:local-volatility:forward}{forward smile} when it does not. In the example, moving from \omterm{def:dv:local-volatility:model}{local volatility} to Heston lowers the one-touch at 7.20 from 43.3\% to 38.9\%.
Stochastic-local volatility with $\lambda=0.5$ sits between them, at 42.2\%.

\section{An FX exotic book}

A currency exotics desk carries thousands of touches, barriers, TARFs and \omterm{def:dv:fx-derivatives:accumulator}{accumulators} across dozens of pairs. Its daily questions follow
from this chapter.
\begin{itemize}
\item \textbf{Which model, which mixing.} One stochastic-local model per pair, its mixing weight calibrated to the liquid touch market, and
the vanna--volga price kept alongside as the market's reference for quoting.
\item \textbf{Where the risk sits.} \omterm{def:dv:greeks-and-the-hedging-pnl:greeks}{Vega} by expiry, and \omterm{def:dv:greeks-and-the-hedging-pnl:greeks}{vanna} and \omterm{def:dv:greeks-and-the-hedging-pnl:greeks}{volga}, which a barrier book concentrates near its barriers. TARF books add
a large exposure to the rate's drift and to jumps, since a sustained move is exactly what they are short.
\item \textbf{What the model cannot see.} Managed currencies move in steps when policy changes, as the August 2015 fixing change showed.
The 2.8\% in two days was a jump for a model calibrated on a pair that had traded near the weak side of a narrow band. Jump risk and the
policy regime belong in the reserves (chapter 27), not in the volatility.
\end{itemize}

\section{Tutorial: from three quotes to a TARF}\label{tut:dv:fx-derivatives:tut}

\begin{tutorial}
\textbf{Goal.} Read the broker quotes from a known market, build the vanna--volga smile and one-touch, calibrate a \omterm{def:dv:fx-derivatives:slv}{stochastic-local volatility
model} by the particle method, and price a TARF. \textbf{End state:} the four figures and the numbers of the weekend problem.
\begin{tutsteps}
\item \textbf{Vanna--volga weights and price}:
\omcode{code/firm/fxvol/firm_fxvol.py}{30}{44}{Vanna--volga weights and the vanilla price.}\label{lst:dv:fx-derivatives:vv}
\item \textbf{The particle method}: bin the particles in spot, estimate the conditional variance, set the leverage and step:
\omcode{code/firm/fxvol/firm_fxvol.py}{108}{139}{Leverage function by the particle method.}\label{lst:dv:fx-derivatives:particle}
\item \textbf{The TARF on simulated fixings}:
\omcode{code/firm/fxvol/firm_fxvol.py}{171}{192}{Target-redemption forward cash flows.}\label{lst:dv:fx-derivatives:tarf}
\item \textbf{Run} \texttt{dv\_fx.quotes()}, \texttt{vv\_smile()}, \texttt{one\_touches()}, \texttt{tarf\_summary()}, \texttt{tarf\_scenario(0.10)} and
\texttt{fig\_fx.py}. The market's \omterm{def:dv:local-volatility:model}{local volatility} is Dupire on a mesh of Heston prices (one Fourier evaluation per expiry).
\end{tutsteps}
\textbf{What to change next.} Fit the mixing weight so that the model's one-touch at 7.60 matches the vanna--volga price; price the TARF under \omterm{def:dv:local-volatility:model}{local
volatility} and under Heston and compare the zero-cost strikes; add a knock-out to the TARF.
\end{tutorial}

\section{Build: FX volatility}\label{bld:dv:fx-derivatives:fxvol}

\begin{build}
\bfield{Purpose} The miniature firm's currency-exotics engine: the market convention for touches and barriers, the stochastic-local model the
desk prices with, and the TARF and \omterm{def:dv:fx-derivatives:accumulator}{accumulator} cash flows its structured-products desk sells.
\bfield{Interface} \texttt{vv\_weights}, \texttt{vv\_price}, \texttt{vv\_implied}; \texttt{up\_hit\_probability}, \texttt{one\_touch\_bs},
\texttt{one\_touch\_vv}; \texttt{SLV(v0, kappa, vbar, eta, rho, mix)}; \texttt{calibrate\_leverage(model, local\_vol, s0, rd, rf, t\_end,
\ldots, n, bins)}; \texttt{simulate\_slv(model, leverage, s0, rd, rf, t\_end, n, \ldots)}; \texttt{tarf\_client\_pnl(fixings, strike, target,
leverage, notional)}.
\bfield{Rules} Quotes are converted with Book 2's conventions before any model sees them; the \omterm{def:dv:fx-derivatives:leverage}{leverage function} is recalibrated with every
change of the smile; the mixing weight is a documented, reviewed input; touches from daily paths are shifted to continuous monitoring.
\bfield{Acceptance tests} \texttt{code/firm/fxvol/tests/}: vanna--volga reproduces its pillars and a flat smile; the Black one-touch matches a
simulation; with no volatility of variance the stochastic-local model is \omterm{def:dv:local-volatility:model}{local volatility} and reprices a flat smile; TARF cash flows on hand
paths.
\bfield{Stretch} Double no-touches under SLV; a mixing weight calibrated to touch quotes; \omterm{def:dv:fx-derivatives:accumulator}{accumulators} with daily knock-outs; the TARF's delta and
\omterm{def:dv:greeks-and-the-hedging-pnl:greeks}{vega} by fixing.
\end{build}

\begin{omsources}
\item M.~B.~Garman and S.~W.~Kohlhagen, ``Foreign currency option values'', \emph{Journal of International Money and Finance} 2(3) (1983)
231--237.
\item J.~Guyon and P.~Henry-Labord\`ere, ``The smile calibration problem solved'', SSRN 1885032 (2011).
\item A.~Perederiy, ``Vanna--volga method for normal volatilities'', SSRN 3262146 (2018); F.~Rolloos, ``On the flat but stochastic implied
volatility assumption in the vanna--volga model'', SSRN 4287718 (2022).
\item Bank for International Settlements, \emph{Quarterly Review} (September 2015), ``EME vulnerabilities take centre stage''.
\end{omsources}

\section{Exercises}

\begin{exercise}[$\star$]\label{exo:dv:fx-derivatives:1}
Show that the Garman--Kohlhagen call is Black's formula on the forward $F=Se^{(r_d-r_f)T}$, discounted at $r_d$.
\end{exercise}

\begin{exercise}[$\star$]\label{exo:dv:fx-derivatives:2}
Why is the correlation between the rate and its variance positive for a dollar against an emerging currency, and what does it do to the risk
reversal?
\end{exercise}

\begin{exercise}[$\star$]\label{exo:dv:fx-derivatives:3}
The client's TARF strike is 7.312 while the forward is 6.861. What has the client paid for a rate 6.6\% better than the forward?
\end{exercise}

\begin{exercise}[$\star\star$]\label{exo:dv:fx-derivatives:4}
Explain why the vanna--volga smile is exact at the three pillars and why it can be wrong at 10 delta.
\end{exercise}

\begin{exercise}[$\star\star$]\label{exo:dv:fx-derivatives:5}
\omterm{def:dv:local-volatility:model}{Local volatility} gives the highest one-touch price at 7.20 among the three models. Explain.
\end{exercise}

\begin{exercise}[$\star\star$]\label{exo:dv:fx-derivatives:6}
Show that with $\lambda=0$ the \omterm{def:dv:fx-derivatives:leverage}{leverage function} is $\sigma_{\mathrm{loc}}(t,S)/\sqrt{v(t)}$, with $v(t)$ the deterministic variance path.
\end{exercise}

\begin{exercise}[$\star\star\star$]\label{exo:dv:fx-derivatives:7}
\emph{Coding.} Price the chapter's TARF under \omterm{def:dv:local-volatility:model}{local volatility} and under Heston and compare the zero-cost strikes and the expected number of
fixings.
\end{exercise}

\begin{exercise}[$\star\star\star$]\label{exo:dv:fx-derivatives:8}
\emph{Find the flaw.} ``Our TARF redeems in 98\% of scenarios, so for a client it is almost a free improvement on the forward rate.''
\end{exercise}

\section{Problem: The TARF}

\begin{problem}[{Weekend problem --- how long it lives and what it costs}]\label{pb:dv:fx-derivatives:1}
A client sells one million dollars a month for a year against the chapter's emerging currency through a TARF: twice the notional when the fixing is
above the strike, redemption when the accumulated gain reaches 0.30 per dollar. The dealer prices with stochastic-local volatility (mixing 0.5).

\medskip\noindent\textbf{Part I --- The market.}
\begin{enumerate}
\item Give the one-year at-the-money volatility, the 25-delta risk reversal and butterfly.
\item What does the positive risk reversal say about the market's fears?
\item Give the one-year forward, and explain why it is below the spot.
\item How well does the stochastic-local model reprice the one-year smile?
\item Give the \omterm{def:dv:fx-derivatives:leverage}{leverage function} near the forward at one year.
\end{enumerate}

\medskip\noindent\textbf{Part II --- The TARF.}
\begin{enumerate}[resume]
\item Give the zero-cost strike and its improvement over the forward.
\item Give the share of paths that redeem and the expected number of fixings.
\item What is the client's gain on the typical path, and where does the zero value come from?
\item Give the zero-cost strikes under \omterm{def:dv:local-volatility:model}{local volatility} and under Heston.
\item Why does a larger target lower the strike the dealer can offer?
\end{enumerate}

\medskip\noindent\textbf{Part III --- The move.}
\begin{enumerate}[resume]
\item The rate rises 10\% over the first quarter. Give the client's expected loss, and how much of it is incurred in the first quarter.
\item Give the same for a 5\% rise.
\item How many fixings does the TARF live, on average, after a 10\% rise?
\item What does the dealer's hedge look like as the accumulated gain approaches the target?
\item Relate the scenario to the renminbi's move in August 2015.
\end{enumerate}

\medskip\noindent\textbf{Part IV --- Judgement.}
\begin{enumerate}[resume]
\item What should the client have been shown before trading?
\item Which model risk matters most for the dealer?
\item Would you sell this product to a company with no natural dollar income?
\item State the \emph{named result}: the expected number of fixings before redemption, and the client's expected loss per million when the rate moves
10\% against it in the first quarter.
\item In one sentence: what has a TARF client sold?
\end{enumerate}
\end{problem}

\section{Interview questions}

\begin{interviewq}[$\star$ \iqroles{trader}]\label{iq:dv:fx-derivatives:1}
Explain the at-the-money, risk reversal and butterfly quotes, and how you would turn them into a smile.
\end{interviewq}

\begin{interviewq}[$\star\star$ \iqroles{researcher}]\label{iq:dv:fx-derivatives:2}
Derive the vanna--volga price of a vanilla. What does the method assume?
\end{interviewq}

\begin{interviewq}[$\star\star$ \iqroles{researcher, developer}]\label{iq:dv:fx-derivatives:3}
What is a \omterm{def:dv:fx-derivatives:slv}{stochastic-local volatility model}, and how would you calibrate its \omterm{def:dv:fx-derivatives:leverage}{leverage function}?
\end{interviewq}

\begin{interviewq}[$\star\star$ \iqroles{trader, risk}]\label{iq:dv:fx-derivatives:4}
Two models that fit the same smile disagree on a one-touch by four points. How do you choose the price?
\end{interviewq}

\begin{interviewq}[$\star\star$ \iqroles{risk, bank}]\label{iq:dv:fx-derivatives:5}
Describe a TARF from the client's side. What can go wrong for the client and for the bank?
\end{interviewq}

\begin{interviewq}[$\star\star\star$ \iqroles{trader, risk}]\label{iq:dv:fx-derivatives:6}
A managed currency may be devalued in a step. How does that change your barrier and TARF book's pricing and hedging?
\end{interviewq}
